\chapter{Cơ sở lý thuyết}
\label{ch:cosolythuyet}

\section{Tính riêng tư vi phân (Differential Privacy)}

Tính riêng tư vi phân (DP)~\cite{dwork2006calibrating} là tiêu chuẩn de-facto cho
riêng tư trong phân tích dữ liệu.

\begin{dinhnghia}[$\eps$-DP]
Cơ chế ngẫu nhiên $M$ thỏa $\eps$-differential privacy nếu với mọi cặp tập dữ
liệu $D_1, D_2$ khác nhau đúng một bản ghi và mọi tập kết quả
$S \subseteq \mathrm{Range}(M)$:
\[
\Pr[M(D_1) \in S] \le e^{\eps}\,\Pr[M(D_2) \in S].
\]
\end{dinhnghia}

Hai tính chất được dùng xuyên suốt luận văn:
\emph{kết hợp tuần tự} (áp dụng $k$ cơ chế $\eps_i$-DP lên cùng dữ liệu cho
$(\sum_i \eps_i)$-DP, kể cả khi cơ chế sau được chọn thích nghi theo đầu ra của
cơ chế trước) và \emph{bất biến hậu xử lý} (mọi phép biến đổi trên đầu ra
\emph{không truy cập dữ liệu gốc} bảo toàn mức riêng tư).

\section{Bất khả phân biệt vị trí địa lý (Geo-Indistinguishability)}

Geo-I~\cite{andres2013geo} tổng quát hóa DP sang không gian metric.

\begin{dinhnghia}[$\eps$-Geo-I]
\label{def:geoi}
Cơ chế $M$ thỏa $\eps$-Geo-Indistinguishability nếu với mọi vị trí thật
$x, x' \in \mathcal{X}$ và mọi tập đầu ra $S$:
\[
\Pr[M(x) \in S] \le e^{\eps\,\dEu(x,x')}\,\Pr[M(x') \in S],
\]
với $\dEu$ là khoảng cách Euclid. Tham số $\eps$ có đơn vị nghịch đảo khoảng
cách (mỗi mét): trong bán kính $r$ bất kỳ, mức phân biệt tối đa là $\eps r$.
\end{dinhnghia}

\paragraph{Cơ chế Laplace phẳng.} Cơ chế kinh điển đạt $\eps$-Geo-I là cộng nhiễu
lấy từ phân phối Laplace phẳng (polar Laplace): góc
$\theta \sim \mathrm{Uniform}[0, 2\pi)$ và bán kính $r$ có mật độ
$p(r) = \eps^2 r e^{-\eps r}$, tức $r \sim \mathrm{Gamma}(2, 1/\eps)$ với kỳ vọng
$\mathbb{E}[r] = 2/\eps$.

\begin{nhanxet}
\label{rmk:expbug}
Lấy $r \sim \mathrm{Exponential}(1/\eps)$ (như hiện thực Thực tập 2) \emph{không}
tạo ra phân phối Laplace phẳng: mật độ hai chiều tương ứng tỷ lệ với
$e^{-\eps r}/r$, phân kỳ khi $r \to 0$, và tỷ số mật độ giữa hai vị trí thật
không bị chặn --- cơ chế thu được không thỏa Định nghĩa~\ref{def:geoi}.
\end{nhanxet}

\section{Cơ chế lũy thừa (Exponential Mechanism)}
\label{sec:em}

Với miền đầu ra rời rạc $V$ và hàm chất lượng $q(x, v)$, cơ chế lũy
thừa~\cite{mcsherry2007mechanism} chọn $v$ với xác suất tỷ lệ
$\exp\!\big(\eps\, q(x,v) / (2\Delta q)\big)$. Trong không gian metric, chọn
$q(x,v) = -\dEu(x,v)$ cho một cơ chế thỏa $\eps$-Geo-I trên $V$ (chứng minh chi
tiết ở Định lý~\ref{thm:rem}); hệ số $1/2$ hấp thụ sự thay đổi của hằng số chuẩn
hóa giữa hai vị trí thật. Đây là nền tảng của cơ chế đề xuất trong
Chương~\ref{ch:phuongphap}.

\section{Mô hình đe dọa}
\label{sec:threatmodel}

Chúng tôi xét kẻ tấn công theo nguyên lý Kerckhoffs: biết \emph{cơ chế} và tham
số của nó, giữ một \emph{prior} trên các vị trí khả dĩ, và quan sát toàn bộ chuỗi
điểm công bố $z_1, \dots, z_T$. Hai hiện thân được dùng để đánh giá
(Chương~\ref{ch:thucnghiem}): tấn công suy luận Bayes tối ưu trên từng điểm, và
tấn công theo dõi HMM khai thác tương quan thời gian giữa các điểm liên tiếp
--- chuẩn đánh giá được xác lập bởi~\cite{shokri2011quantifying,xiao2015protecting}.
